\documentclass[11pt]{article}

\usepackage[margin=1.1in]{geometry}
\usepackage{amsmath,amssymb,amsthm,mathtools}
\usepackage{enumitem}
\usepackage{xcolor}
\usepackage[colorlinks=true,linkcolor=blue!55!black,citecolor=blue!55!black,urlcolor=blue!55!black]{hyperref}

\newcommand{\F}{\mathbb{F}}
\newcommand{\eq}{\mathrm{eq}}
\newcommand{\cube}[1]{\{0,1\}^{#1}}
\newcommand{\pol}[1]{\widehat{#1}}
\newcommand{\inner}[2]{\langle #1,\, #2 \rangle}
\newcommand{\val}[1]{\langle #1 \rangle}
\newcommand{\maj}{\mathrm{maj}}

\emergencystretch=1.5em

\theoremstyle{plain}
\newtheorem{lemma}{Lemma}[section]

\title{\textbf{The jagged adapter over a weighted-claim PCS}}
\author{}
\date{August 3, 2026}

\begin{document}
\maketitle

\begin{center}
\fbox{\textbf{Warning: AI-assisted writing, some parts have not been properly reviewed yet.}}
\end{center}

%=====================================================================
\section{The goal}\label{sec:goal}
%=====================================================================

An arithmetization presents $N$ columns; here and below $N,m,n,k,h,b,B,L,t,e,A$ and $\beta$
are nonnegative integers. Column $i$ is a multilinear
$P_i:\cube{k_i}\to\F$, so of logical length $2^{k_i}$, but only a prefix of $h_i\le2^{k_i}$
entries carries data: from $h_i$ upward the column repeats a fixed public value $p_i\in\F$.
The heights $h_i$ are run-time quantities while the dimensions $k_i$ are fixed when the
circuit is, so the padding is real waste, and committing to the columns one at a time costs
the verifier work linear in $N$ on top of it.

Commit instead, once, to the concatenation of the real prefixes with the padding left out: a
single dense $q:\cube m\to\F$. Each column then occupies a known interval of the dense cube. What has to be built is a reduction turning the
arithmetization's claims, which are evaluations $\pol{P_i}(\zeta)$ of the padded columns, into
one weighted claim $\inner Wq=C$ on that single commitment, against a weight $W$ the verifier
can evaluate for itself. This is the jagged PCS of~\cite{jagged}.

Two things distinguish what follows from that paper. Our PCS opens an arbitrary
verifier-evaluable weighted claim, so the paper's $m$-round reduction sumcheck, about
$5\cdot2^m$ prover multiplications, is unnecessary and $W$ is opened directly. And it runs its
own sumcheck over all $m$ coordinates, so it asks for $\pol W$ at one point of $\F^m$ and for
nothing else.

Notation, multilinear extensions and the inner-product commitment interface are those of the
main document: an opening proves $\inner Wq=\sum_{w\in\cube m}W(w)q(w)=C$ against any weight
$W$ the verifier can evaluate, a point evaluation being $W=\eq(\zeta,\cdot)$. Bit strings are
little-endian, with $\val a=\sum_j a_j2^j$.

%=====================================================================
\section{The layout}\label{sec:layout}
%=====================================================================

\paragraph{Blocks} A \emph{block} is a set of $B=2^b$ columns sharing a logical dimension $k$, a
real height $h\le2^k$, and an opening point: whenever one of them is claimed, all of them are,
at the same point. A set of mutually compatible columns is cut into blocks of power-of-two
widths greedily, nine columns becoming $8+1$, so no dummy column is ever introduced. Write
$c\in\cube b$ for a column's \emph{slot} in its block and $P_c$ for that column.

\paragraph{Placement} A block is stored row-major at an integer offset $t\ge0$:
\begin{equation}
  q\bigl(t+\val c+B\val z\bigr)=P_c(z)-p_c,\qquad c\in\cube b,\ z\in\cube k,\ \val z<h,
  \label{eq:place}
\end{equation}
so it occupies the integer interval $[t,e)$ with $e=t+Bh$. What is stored is each column with
its own pad value subtracted. Since the column equals the constant $p_c$ at and above $h$, the
shifted column vanishes there, so it is supported on the committed prefix.

\paragraph{The dense cube.} Blocks are concatenated with no alignment gap, and only the last
is followed by zero padding. The cube dimension $m$ is the least one admitting both the packed
area and every block's own logical cube,
\[
  2^m\ge A:=\sum_{\text{blocks}}Bh,\qquad m\ge b+k\ \text{ for every block},
\]
the second condition being what lets a block's $b+k$ coordinates be read inside $\cube m$.
Every endpoint is public and none is transmitted: they are the prefix sums of the block
lengths $Bh$, which a circuit verifier forms with a ripple-carry adder.

%=====================================================================
\section{The reduction}\label{sec:reduction}
%=====================================================================

Fix a block, with its interval $[t,e)$, width $B=2^b$, dimension $k$, height $h$, pad values
$p_c$, and opening point $\zeta\in\F^k$. Its input is one claim per slot,
$\pol{P_c}(\zeta)=v_c$ with $v_c\in\F$ and $c\in\cube b$.

\paragraph{The padding} Set $R_c:=P_c-p_c$, which by~\eqref{eq:place} is what the commitment holds
on $[0,h)$ and which vanishes at and above $h$. The multilinear extension of a constant is
that constant, so $\pol{P_c}(\zeta)=\pol{R_c}(\zeta)+p_c$ at every $\zeta$, whatever $h$ is.
Correcting each claimed value, $v_c':=v_c-p_c$, the input becomes $\pol{R_c}(\zeta)=v_c'$.

\paragraph{The weight} For integer bounds $0\le t\le e\le2^m$ and a family of \emph{row weights}
$u=(u_j^0,u_j^1)_{j<n}\in(\F^2)^n$, define the \emph{interval weight} $W^{t,e}_u:\cube m\to\F$
by
\begin{equation}
  W^{t,e}_u(w):=\sum_{\substack{a\in\cube n\\ t+\val a<e}}\Bigl(\prod_{j<n}u_j^{a_j}\Bigr)\bigl[\val w=t+\val a\bigr].
  \label{eq:weight}
\end{equation}

\paragraph{The block} The $B$ claims share $\zeta$ but not their values, so combining them needs
one coefficient per claim. Taking those coefficients to be \emph{consecutive} powers of one challenge
$\gamma\in\F$ is what makes the combination collapse to a single interval weight. Put
\begin{equation}
  u^{\star}(\gamma,\zeta)_j:=\bigl(1,\,\gamma^{2^j}\bigr)\ \ (j<b),\qquad u^{\star}(\gamma,\zeta)_{b+i}:=(1-\zeta_i,\,\zeta_i)\ \ (i<k),
  \label{eq:ustar}
\end{equation}
a family in $(\F^2)^{b+k}$, so $n=b+k$. Expanding each $\pol{R_c}$ over its support and using
$\val c+B\val z=\val{(c,z)}$,
\begin{equation}
  \sum_{c\in\cube b}\gamma^{\val c}\,\pol{R_c}(\zeta)
  =\sum_{c\in\cube b}\ \sum_{\val z<h}\gamma^{\val c}\,\eq(\zeta,z)\,q\bigl(t+\val{(c,z)}\bigr)
  =\bigl\langle W^{t,e}_{u^{\star}(\gamma,\zeta)},\,q\bigr\rangle .
  \label{eq:batch}
\end{equation}
The last step is~\eqref{eq:weight} read at $a=(c,z)$: the selector factors of $u^{\star}$
contribute $\prod_{j<b}(\gamma^{2^j})^{c_j}=\gamma^{\val c}$, the remaining factors contribute
$\eq(\zeta,z)$, and the condition $t+\val a<e=t+Bh$ is exactly $\val z<h$ because $\val c<B$.

\noindent\fbox{\parbox{0.955\textwidth}{
\emph{Public:} the blocks, each with its interval $[t,e)$, width $B=2^b$, dimension $k$,
height $h$, and pad values $p_c\in\F$.

\smallskip
\textbf{Commit.} The prover commits to $q$.

\smallskip
\textbf{Open.} Each block is claimed at its own point $\zeta\in\F^k$, one claim per column:
$\pol{P_c}(\zeta)=v_c\in\F$ for $c\in\cube b$. Below, $t$, $e$, $B$, $b$, $k$, $h$, $p_c$,
$\zeta$ and $\beta$ denote the data of whichever block a term belongs to.
\begin{enumerate}[itemsep=1pt,topsep=2pt]
\item Both parties correct each value for the padding: $v_c':=v_c-p_c$.
\item The verifier observes every $v_c'$, then samples $\gamma\leftarrow\F$ and ranks the
claims. Fix any public order of the blocks and number the claims $0,1,2,\dots$ along it, a
block's $B$ claims consecutively and in slot order. Write $\beta$ for the number its slot
$0^b$ receives, so slot $c$ has rank $\beta+\val c$.
\item Both parties form
\[
  W:=\sum_{\text{blocks}}\gamma^{\beta}\,W^{t,e}_{u^{\star}(\gamma,\zeta)},
  \qquad C:=\sum_{c}\gamma^{\mathrm{rank}(c)}\,v_c' .
\]
\item The PCS opens $\inner Wq=C$.
\end{enumerate}
}}

\medskip
Multiplying~\eqref{eq:batch} by $\gamma^{\beta}$ and summing over blocks gives
$\inner Wq=\sum_c\gamma^{\mathrm{rank}(c)}\pol{R_c}(\zeta)$, one term per claim, so
\[
  \inner Wq-C=\sum_{c}\gamma^{\mathrm{rank}(c)}\bigl(\pol{R_c}(\zeta)-v_c'\bigr).
\]
If every claim holds this vanishes. If one fails it is a nonzero polynomial in $\gamma$ of
degree below the number of claims, so it vanishes for at most that many challenges. Past that
the soundness is the PCS's own. Nothing else about the reduction needs proving; it is an
identity. One task is left for the verifier, evaluating $\pol W$, and the rest of the note is
about it.

%=====================================================================
\section{The weight is a branching program}\label{sec:automaton}
%=====================================================================

By~\eqref{eq:weight} the weight tests an integer addition and an inequality, bit by bit. Two
registers suffice for that, which is why the automaton below has width four~\cite{jagged}.

Fix $t,e,u$ as in~\eqref{eq:weight} with $t<e$ and $n\le m$, and an evaluation point
$x\in\F^m$. An empty interval needs no automaton, its weight being identically zero. The
automaton reads $m+1$ positions $j=0,\dots,m$, least significant first: one per bit of a dense
index, plus a last one that rejects an addition carrying out of the cube and lets an interval
ending at $2^m$ be represented. Positions past the data are filled by convention, $u_j:=(1,0)\in\F^2$ for
$n\le j\le m$ and $x_m:=0$.

The automaton keeps two bits. The first, $\kappa$, is the carry of the addition
$\val w=t+\val a$ over the bits read so far. The second, $\lambda$, records whether
$\val w<e$ over those bits; a more significant differing bit overrides a less significant one,
so reading upward each differing bit simply overwrites $\lambda$. Index the four pairs
$(\kappa,\lambda)$ by the integer $\kappa+2\lambda$, so the state set is $Q:=\{0,1,2,3\}$.

What the automaton propagates is not one state but a mass per state, a vector in $\F^Q$. It
starts at the indicator of state $0$, no carry and nothing yet smaller, and it accepts in state
$2$, no carry and $\val w<e$ settled.

At position $j$ let $t_j,e_j\in\{0,1\}$ be the $j$-th bits of $t$ and $e$, and write
$\chi_j^0:=1-x_j$ and $\chi_j^1:=x_j$ in $\F$ for the two per-coordinate factors of $\eq$, so
that $\eq(\varepsilon,x)=\prod_j\chi_j^{\varepsilon_j}$ for Boolean $\varepsilon$. Set
\[
  g_0:=u_j^0\chi_j^0,\quad g_1:=u_j^1\chi_j^0,\quad g_2:=u_j^0\chi_j^1,\quad g_3:=u_j^1\chi_j^1 .
\]
Reading a row bit $a$ and an index bit $\omega$ from $(\kappa,\lambda)$ is permitted only when
$\omega=a\oplus\kappa\oplus t_j$, and then
\[
  \kappa'=\maj(a,\kappa,t_j),\qquad
  \lambda'=\begin{cases}\lambda & \omega=e_j,\\ e_j & \omega\neq e_j,\end{cases}
  \qquad\text{with weight } u_j^a\chi_j^{\omega}.
\]

\begin{lemma}\label{lem:automaton}
For every $x\in\F^m$ the automaton outputs $\pol{W^{t,e}_u}(x)$, namely
\[
  \sum_{\substack{a\in\cube n\\ t+\val a<e}}\Bigl(\prod_{j<n}u_j^{a_j}\Bigr)\eq\bigl(t+\val a,\,x\bigr).
\]
\end{lemma}

\begin{proof}
Expanding the product of transition matrices, a path is a choice of $(a_j,\omega_j)_{j\le m}$
of weight $\prod_j u_j^{a_j}\chi_j^{\omega_j}$. The constraint
$\omega_j=a_j\oplus\kappa_j\oplus t_j$ with $\kappa_{j+1}=\maj(a_j,\kappa_j,t_j)$ and
$\kappa_0=0$ is binary addition, so paths biject with choices of $a$ and satisfy
$\val\omega=t+\val a$; then $\prod_j\chi_j^{\omega_j}=\eq(\omega,x)$, and $u_j=(1,0)$ for
$j\ge n$ kills every $a$ with a bit at or above $n$. Acceptance is $t+\val a<e$: $\lambda$
ends at $1$ exactly when the most significant position where $\omega$ and $e$ differ has
$e$-bit $1$, and a path that carries into position $m$ is annihilated there, because the carry
forces $\omega_m=1$, of weight $\chi_m^1=x_m=0$. When the interval ends at $2^m$ that same
position carries $e_m=1$, which sets $\lambda$ to $1$. Both sides are multilinear in every coordinate of $x$ and
agree at Boolean arguments, hence agree everywhere.
\end{proof}

Since $t_j$ and $e_j$ are public bits, a step selects one of four fixed matrices rather than
evaluating a tensor over $(a,\omega)$. With $S$ incoming and $S'$ outgoing:
\[
\renewcommand{\arraystretch}{1.25}
\begin{array}{l@{\qquad}l}
(t_j,e_j)=(0,0): &
\begin{aligned}
S_0'&=S_0(g_0+g_3)+(S_1+S_3)g_2+S_2g_3, & S_1'&=S_1g_1,\\
S_2'&=S_2g_0, & S_3'&=S_3g_1,
\end{aligned}\\[2.2ex]
(t_j,e_j)=(0,1): &
\begin{aligned}
S_0'&=S_0g_3+S_1g_2, & S_1'&=0,\\
S_2'&=S_0g_0+S_2(g_0+g_3)+S_3g_2, & S_3'&=(S_1+S_3)g_1,
\end{aligned}\\[2.2ex]
(t_j,e_j)=(1,0): &
\begin{aligned}
S_0'&=(S_0+S_2)g_2, & S_1'&=S_0g_1+S_1(g_0+g_3)+S_3g_3,\\
S_2'&=0, & S_3'&=S_2g_1+S_3g_0,
\end{aligned}\\[2.2ex]
(t_j,e_j)=(1,1): &
\begin{aligned}
S_0'&=S_0g_2, & S_1'&=S_1g_3,\\
S_2'&=S_2g_2, & S_3'&=(S_0+S_2)g_1+S_1g_0+S_3(g_0+g_3).
\end{aligned}
\end{array}
\]
Each arm costs $6$ multiplications and the $g_i$ cost one more. At a row coordinate, where
$u_j=(1-r,r)$ for that coordinate's value $r\in\F$, form $\pi:=rx_j$ and read off
$g_0=1-r-x_j+\pi$, $g_1=r-\pi$, $g_2=x_j-\pi$, $g_3=\pi$. At a selector coordinate, where
$u_j=(1,\gamma^{2^j})$, form $\pi:=\gamma^{2^j}x_j$ and read off $g_0=1-x_j$,
$g_1=\gamma^{2^j}-\pi$, $g_2=x_j$, $g_3=\pi$. So a position costs $7$ multiplications, plus
one squaring per selector coordinate to advance $\gamma^{2^j}$.

The last position is free. There $u_m=(1,0)$ and $x_m=0$, hence $g_0=1$ and
$g_1=g_2=g_3=0$, and $t_m=0$ because $t<2^m$, so the arms above leave
\begin{equation}
  \pol{W^{t,e}_u}(x)=S_2+e_mS_0
  \label{eq:accept}
\end{equation}
in terms of the state $S$ after position $m-1$: the mass with no pending carry and $\val w<e$
already decided, plus, only when the interval ends exactly at $2^m$, the mass still tied at
equality. A full run therefore costs $7m$.

One run serves a whole block, which is what~\eqref{eq:batch} bought, and it is why the weight
was defined for row pairs that need not sum to $1$: normalising $u^{\star}$, that is writing
$(1,\gamma^{2^j})=(1+\gamma^{2^j})\cdot(1-\theta_j,\,\theta_j)$ with
$\theta_j=\gamma^{2^j}/(1+\gamma^{2^j})$, would cost $b$ inversions and fail wherever
$1+\gamma^{2^j}$ vanishes.

%=====================================================================
\section{What the verifier evaluates}\label{sec:evaluate}
%=====================================================================

The PCS is BaseFold/WHIR-style~\cite{ligerito,whir}: it consumes $m-L$ of the $m$ variables in
sumcheck rounds and then transmits the remaining $L$-variable message $f:\cube L\to\F$ in the
clear, with $L$ small, typically $3$ or $4$. That message is public, so the verifier evaluates
it itself and the sumcheck runs on over its $L$ coordinates. The opening therefore ends at one
point $x\in\F^m$, the $m-L$ round challenges followed by the $L$ tail challenges, and the
verifier's whole remaining task is
\[
  \pol W(x)=\sum_{\text{blocks}}\gamma^{\beta}\,\pol\Omega(x),
  \qquad \Omega:=W^{t,e}_{u^{\star}(\gamma,\zeta)},
\]
one automaton run per block.

\medskip
\noindent\fbox{\parbox{0.955\textwidth}{
\textbf{The verifier's work for one block.}
\emph{Input:} the bits of $t$ and $e$; the log-width $b$; the scalar $\gamma$ and the number
$\beta$ of the block's slot $0^b$; the row point $\zeta\in\F^k$; the point $x\in\F^m$.
\emph{Output:} the block's contribution $\gamma^{\beta}\pol\Omega(x)$.
\begin{enumerate}[itemsep=1pt,topsep=2pt]
\item Run the automaton over the positions $j=0,\dots,m-1$ at coordinates $x_j$, with
$u_j=(1,\gamma^{2^j})$ for $j<b$, then $u_j=(1-\zeta_{j-b},\,\zeta_{j-b})$ for
$b\le j<b+k$, then $u_j=(1,0)$, giving $S\in\F^Q$.
\item Return $\gamma^{\beta}(S_2+e_mS_0)$ by~\eqref{eq:accept}.
\end{enumerate}
Cost $7m$ multiplications plus $b$ squarings.
}}

%=====================================================================
\section{Cost}\label{sec:cost}
%=====================================================================

The prover makes one commitment, to $q$. Building $q$ itself is $O(A)$, a row-major copy with
one subtraction per cell. Its opening work is building $W$ over $\cube m$: one $\eq$ table per
distinct opening point, then one multiply-accumulate per (cell, claim) pair into the cells of
that claim's column.

Given the endpoint bits, the powers of $\gamma$ and $C$, all of which are $O(m)$ or one pass
over the claims, the verifier spends on the weight itself $7m$ multiplications plus $b$
squarings per \emph{block} rather than per column. That computation is straight-line in $m$ and
the number of blocks, with the shape of the data entering only as bits, which is what makes it
cheap to prove recursively.

%=====================================================================
\begin{thebibliography}{9}

\bibitem{jagged} Tamir Hemo, Kevin Jue, Eugene Rabinovich, Gyumin Roh, and Ron D. Rothblum.
Jagged Polynomial Commitments (or: How to Stack Multilinears). Manuscript, April 2026.

\bibitem{ligerito} Andrija Novakovic and Guillermo Angeris. Ligerito: A Small and Concretely
Fast Polynomial Commitment Scheme. Cryptology ePrint Archive, Paper 2025/1187, 2025.

\bibitem{whir} Gal Arnon, Alessandro Chiesa, Giacomo Fenzi, and Eylon Yogev. WHIR:
Reed--Solomon Proximity Testing with Super-Fast Verification. In \emph{EUROCRYPT 2025}.

\end{thebibliography}

\end{document}
